% ===================================================================
%  TABELL Nr 7 — Byn om vintern. — Lanthuset om våren.
%  Traduction 2026 d'apres texto/io, controlee sur texto/fr et
%  texto/es. Consigne : voir texto/sv/10-tabell-01.tex.
% ===================================================================

%%K t07-noto-1 noto
\VUnotes{143.2mm}{%
(*) Se om måltidernas enskildheter, tabellerna 6 och 13.%
}

%%K t07-apar-1 apar
\VUsaut{4.0mm}
\VUcentre{13.2pt}{90}{ANDRA SERIEN}
\VUsaut{3.0mm}
\VUfilet{16mm}
\VUsaut{5.6mm}
\VUtitre{15.0pt}{60}{TABELL Nr 7}
\VUsaut{3.0mm}

%%K t07-tit-1 sub
\VUcentre{12.0pt}{0}{\textit{Första scenen.}}
\VUsaut{3.0mm}

%%K t07-tit-2 sub
\VUpk{10.2pt}{40}{Byn om vintern.}
\VUsaut{4.0mm}

%%K t07-c1-01-1 p
1. --- Under nyårsferierna följde jag min far till Nystad, en liten by där han måste uträtta
några affärer.

%%K t07-c1-02-1 p
2. --- Vi tog \VUgras{diligensen} \textsuperscript{(11)} som går mellan staden och den
köpingen; men eftersom det var mycket kallt, var den resan i ett föga bekvämt åkdon inte
särskilt behaglig.

%%K t07-c1-03-1 p
3. --- Därför kände vi ett sant nöje när vi såg \VUgras{kusken} stanna sin vagn på torget,
framför \textit{Vita Björnens} \VUgras{värdshus} \textsuperscript{(2)}.
\VUgras{Värdshusvärden} \textsuperscript{(4)} väntade på oss på tröskeln till sitt hus, medan
han lugnt rökte sin \VUgras{pipa}.

%%K t07-c1-03-2 p
Han bjöd oss stiga in och jag lät mig inte bedjas att slå mig ner framför rishögens eld som
sprakade i härden. Då hånade jag den skarpa
\VUgras{nordanvinden} som får den gamla \VUgras{skylten} \textsuperscript{(3)}
att gnissla och som i svarta virvlar förde bort \VUgras{röken} \textsuperscript{(1)} som steg
ur skorstenen.

%%K t07-c1-04-1 p
4. --- Det var frukostdags. Utan att vänta längre åt vi med god aptit den enkla men
välsmakande måltid som man serverade oss (*).

%%K t07-tit-3 sub
\VUpk{10.2pt}{40}{Några besök.}
\VUsaut{3.0mm}

%%K t07-c1-05-1 p
5. --- Efter frukosten följde jag min far, som är försäkringsagent, till hans kunder. Först
besökte vi \VUgras{smeden} \textsuperscript{(5)}, som bor bredvid värdshuset. Han höll på att
sko en häst. Jag beundrade styrkan hos hans kamrat som stadigt höll hästens \VUgras{hov} på
sitt läderförkläde, medan smeden fäste
\VUgras{skon} \textsuperscript{(6)} med kraftiga hammarslag.

%%K t07-c1-06-1 p
6. --- Sedan gick vi till byns \VUgras{skomakare} \textsuperscript{(7)}, gamle Jakob. Fastän
han har en praktfull \VUgras{stövel} som skylt, nöjde han sig med att halvsula ett gammalt par
genomslitna skor.

%%K t07-c1-06-2 p
Den gamle sade skrattande till oss: « I staden var jag “skomakare” och jag hade inga kunder.
Här är jag “skolagare” och arbetet är mycket. »

%%K t07-c1-07-1 p
7. --- Han talade med oss om sin granne \VUgras{barberaren} \textsuperscript{(8)}, vars kunder
inte är nöjda. Hans \VUgras{rakknivar} är alltid hackiga och hans \VUgras{saxar} klipper
aldrig bra. Men eftersom han är byns ende barberare, är man naturligtvis tvungen att låta
honom raka en och klippa ens hår. För övrigt är han en girig och ovänlig man.

%%K t07-c1-08-1 p
8. --- Lyckligtvis liknar de övriga \VUgras{grannarna} honom inte. Till exempel är
\VUgras{vagnmakaren} \textsuperscript{(9)} en godhjärtad man, flitig och hjälpsam. Det är ett
nöje att låta honom laga en vagn. Hans arbete är alltid avslutat.

%%K t07-c1-09-1 p
9. --- Gamle Jakob klagade inte heller på \VUgras{sadelmakaren} \textsuperscript{(10)}, som
han ibland hjälper att sy \VUgras{seldonen} när arbetet brådskar.

%%K t07-c1-09-2 p
Min far frågade honom om hans familj: « Alla mår bra. Mitt barnbarn går i
\VUgras{skolan} \textsuperscript{(12)}. Hennes \VUgras{lärarinna}
\textsuperscript{(13)} är mycket nöjd med henne, hon är en god \VUgras{elev}
\textsuperscript{(14)}. Hon liknar inte min odåga till son; han tänker bara på att leka.
\VUgras{Läraren} \textsuperscript{(15)} lyckas inte lära honom någonting alls. »

%%K t07-tit-4 sub
\VUsaut{3.0mm}
\VUpk{10.2pt}{40}{Byskvaller.}
\VUsaut{3.0mm}

%%K t07-c1-10-1 p
10. --- Den gamle berättade sedan byns nyheter för oss. Man ska bygga en ny ska, för den som
är inrymd i \VUgras{kommunalhuset} har blivit för liten. För övrigt är det lätt, för det ska
räcka att köpa en del av \VUgras{mjölnarens} trädgård. Det ska bli mycket lönsamt för den
stackars mannen. På grund av kylan kan \VUgras{kvarnens} \textsuperscript{(16)}
\VUgras{kvarnstenar} inte längre mala säden, för \VUgras{hjulet} \textsuperscript{(17)} har
frusit fast av
\VUgras{isen} \textsuperscript{(25)} i \VUgras{bäcken} \textsuperscript{(23)}.

%%K t07-c1-11-1 p
11. --- « Vad byggnader beträffar, har ni sett vår nya \VUgras{kyrka} \textsuperscript{(18)}?
Den är mycket pittoresk under sin snömantel.
\VUgras{Kråkorna} \textsuperscript{(19)}, som inte längre känner igen sitt gamla
klocktorn, sitter sorgset på det stora trädet på \VUgras{kyrkogården} \textsuperscript{(20)}
».

%%K t07-c1-12-1 p
12. --- Den tanken på kyrkogården påminde skomakaren om de personer som min far hade känt och
som dog förra året. « Hur många \VUgras{gravar} \textsuperscript{(21)}, min gode herre, har
inte lagts till de andra, under
\VUgras{cypressen} \textsuperscript{(22)}! Döden har mejat ner vår gamle notarie.
Alla byborna bevistade hans \VUgras{begravning}. »

%%K t07-c1-13-1 p
13. --- « Här är hans efterträdare som åker skridsko på bäcken. Han kom just till mig för att
låta laga remmen på sin \VUgras{skridsko} \textsuperscript{(26)} som hade gått av. »

%%K t07-c1-14-1 p
14. --- « Här är också på bäckens strand, den nye \VUgras{fjärdingsmannen}
\textsuperscript{(27)}. Han är en före detta gendarm som skrämmer byns odågor. De flyr så
snart de får syn på hans \VUgras{damasker} \textsuperscript{(29)} och hans \VUgras{kaskett}
\textsuperscript{(28)}. »

%%K t07-c1-15-1 p
15. --- « Ha! här är gamle Josef som kör en \VUgras{kärra med rishögar} \textsuperscript{(30)}
åt bagaren. --- Det är sant, sade min far, jag känner igen hans spann av \VUgras{mulåsnor}
\textsuperscript{(31)}. Jag behöver verkligen tala med honom, jag ska begagna tillfället.
Farväl, fader Jakob. »

%%K t07-c1-16-1 p
16. --- Just när vi gick ut, lämnade byns barn skolan och rusade springande till
\VUgras{isbanan} \textsuperscript{(33)} med risk att falla och bryta en lem vid
\VUgras{fallet} \textsuperscript{(34)}. En av dem tog sin \VUgras{kälke}
\textsuperscript{(32)} hos vagnmakaren, satte en av sina kamrater på den, och for i väg
springande.

%%K t07-c1-17-1 p
17. --- Andra rusade till en \VUgras{snöhög} \textsuperscript{(37)} bredvid
\VUgras{bron} \textsuperscript{(40)} och började bombardera \VUgras{snöbilden}
\textsuperscript{(39)} som de hade rest dagen förut och på vilken de hade satt en hatt, en
pipa och en skolväska över axeln.

%%K t07-c1-18-1 p
18. --- De bekymrar sig föga om den isande vinden och kylan. De flesta har inte ens någon
\VUgras{tjock halsduk} \textsuperscript{(35)} och, för att värma sig efter striden med
\VUgras{snöbollar} \textsuperscript{(38)}, räcker en handfull mycket varma \VUgras{kastanjer}
\textsuperscript{(42)} åt dem, köpta av den gamla
\VUgras{försäljerskan} Jakobina, som darrar av köld under sin alltför tunna
\VUgras{trasiga sjal} \textsuperscript{(43)}.

%%K t07-tit-5 sub
\VUsaut{3.0mm}
\VUcentre{12.0pt}{0}{\textit{Andra scenen.}}
\VUsaut{3.0mm}

%%K t07-tit-6 sub
\VUpk{10.2pt}{40}{Lanthuset om våren.}
\VUsaut{4.0mm}

%%K t07-c2-01-1 p
1. --- Trots alla vinterns nöjen hälsar vi med djup glädje \VUgras{vårens} återkomst.

%%K t07-c2-01-2 p
Vilken upplivande tavla erbjuder inte en bondgård, om kvällen, i maj månad!

%%K t07-c2-02-1 p
2. --- Vid horisonten går \VUgras{solen} \textsuperscript{(1)} långsamt ner, och översvämmar
\VUgras{landsbygden} \textsuperscript{(3)} med sina purpurfärgade
\VUgras{strålar} \textsuperscript{(2)}. Den flitige \VUgras{plöjaren}
\textsuperscript{(6)} skyndar sig att plöja sin \VUgras{åker} \textsuperscript{(4)}.

%%K t07-c2-02-2 p
Med sin \VUgras{plog} \textsuperscript{(5)} spänd för en kraftig \VUgras{häst}
\textsuperscript{(7)}, drar han \VUgras{fåran} \textsuperscript{(8)} i vilken
\VUgras{såningsmannen} \textsuperscript{(9)} med full hand ska kasta
\VUgras{utsädet} som ska ge de gyllene axen.

%%K t07-c2-03-1 p
3. --- \VUgras{Slåtterkarlarna} \textsuperscript{(11)} med sina liar har slagit ängens gräs,
hövändarna har torkat \VUgras{höet} \textsuperscript{(10)} genom att vända det med
\VUgras{högafflar} \textsuperscript{(12)} och räfsor, och här är det att de kommer tillbaka.

%%K t07-c2-03-2 p
Somliga går framför den tunga vagnen som dras av oxar; de bär sitt redskap över axeln. Andra,
latare, har bekvämt slagit sig ner på det doftande höet.

%%K t07-c2-04-1 p
4. --- I förbifarten hälsar de vänskapligt på \VUgras{trädgårdsmästaren}
\textsuperscript{(16)} som i \VUgras{fruktträdgården} \textsuperscript{(13)} är sysselsatt med
att beskära \VUgras{fruktträden} och ympa \VUgras{päronträden} \textsuperscript{(14)}.
\VUgras{Plommonträden} \textsuperscript{(15)} börjar blomma, och, i den väl gödslade
\VUgras{köksträdgården} \textsuperscript{(17)}, står grönsakerna, kålen och salladen redan
bra.

%%K t07-c2-04-2 p
På den lilla muren breder en vacker \VUgras{påfågel} \textsuperscript{(18)} ut sin stjärt, för
att låta beundra sina praktfulla fjädrar.

%%K t07-tit-7 sub
\VUpk{10.2pt}{40}{Bondgården.}
\VUsaut{3.0mm}

%%K t07-c2-05-1 p
5. --- \VUgras{Stallets} \textsuperscript{(19)} dörrar står öppna och man ser krubban och
\VUgras{halmbädden} \textsuperscript{(23)} åt \VUgras{stoet} \textsuperscript{(21)} som
\VUgras{stalldrängen} \textsuperscript{(20)} just har spänt ifrån. \VUgras{Fölet}
\textsuperscript{(22)} så komiskt med sina långa ben, följer sin mor skuttande.

%%K t07-c2-06-1 p
6. --- Vid \VUgras{brunnen} \textsuperscript{(29)} går \VUgras{korna} \textsuperscript{(25)}
för att dricka ur \VUgras{tråget} \textsuperscript{(26)} av trä som tjänar dem som vattenho.
\VUgras{Kalven} \textsuperscript{(24)} begagnar tillfället för att dia, och \VUgras{tjuren}
\textsuperscript{(27)} som känner den goda doften från höladan, råmar glatt och lyfter stolt
huvudet med de kraftiga \VUgras{hornen} \textsuperscript{(28)}.

%%K t07-c2-07-1 p
7. --- Alla arbetar. \VUgras{Tjänsteflickan} \textsuperscript{(32)} håller brunnens
\VUgras{rep} \textsuperscript{(31)} i handen. Hon öser vatten med en
\VUgras{hink} \textsuperscript{(30)} för att vattna boskapen. Bredvid
\VUgras{ladan} \textsuperscript{(33)} räfsar en man ihop halmen och framför
\VUgras{bostadshusets} \textsuperscript{(34)} dörr spinner en gammal
\VUgras{spinnerska} \textsuperscript{(35)} med sin spinnrock och sin
\VUgras{slända}, \VUgras{linet} eller \VUgras{hampan} som i forna tider.

%%K t07-tit-8 sub
\VUsaut{3.0mm}
\VUpk{10.2pt}{40}{Bonden.}
\VUsaut{3.0mm}

%%K t07-c2-08-1 p
8. --- \VUgras{Bonden} \textsuperscript{(36)} leder alla de arbetena och ger anvisningar åt
sina drängar. Han befaller att ta in under tak \VUgras{harven} \textsuperscript{(40)} och
\VUgras{hackan} \textsuperscript{(39)} som man har glömt bredvid \VUgras{kojan}
\textsuperscript{(38)} där \VUgras{hunden} \textsuperscript{(37)} bor.

%%K t07-c2-09-1 p
9. --- Hans rop skrämmer en \VUgras{duva} \textsuperscript{(41)}, som med alla vingar flyger
in i \VUgras{duvslaget} \textsuperscript{(42)}, och
\VUgras{hönorna} \textsuperscript{(43)} som skyndar sig att gå in i
\VUgras{hönshuset} \textsuperscript{(44)} igen.

%%K t07-c2-10-1 p
10. --- Framför \VUgras{fårhuset} \textsuperscript{(48)}, här är den gamle
\VUgras{herden} \textsuperscript{(45)} som för tillbaka sin \VUgras{(får)hjord}
\textsuperscript{(49)}. Med sin \VUgras{herdestav} \textsuperscript{(46)} i handen, som aldrig
fattas honom, lika litet som hans blekta \VUgras{blus} \textsuperscript{(47)}, driver han till
stallet \VUgras{baggar} \textsuperscript{(50)}, \VUgras{tackor} \textsuperscript{(53)},
\VUgras{lamm} \textsuperscript{(54)}, \VUgras{getter} \textsuperscript{(51)} och
\VUgras{killingar} \textsuperscript{(52)}. Hans trogne \VUgras{hund}
\textsuperscript{(55)} Argus kommer sist och eggar med sitt skällande eftersläntrarnas flit.

%%K t07-c2-11-1 p
11. --- Allt det bullret och den rörelsen stör inte lugnet hos den goda
\VUgras{mjölkkon} \textsuperscript{(56)} som låter sin \VUgras{skälla}
\textsuperscript{(57)} pingla, medan man mjölkar henne framför \VUgras{ladugårdens}
\textsuperscript{(58)} dörr.

%%K t07-c2-12-1 p
12. Hon tycks förstå, att hon måste ge sin mjölk, för att \VUgras{bondhustrun}
\textsuperscript{(60)} skulle kunna göra \VUgras{smör} i \VUgras{kärnan}
\textsuperscript{(61)} eller de utmärkta \VUgras{ostarna} \textsuperscript{(64)}, som droppar
från hyllan som ligger på \VUgras{baljan} \textsuperscript{(63)}.

%%K t07-c2-13-1 p
13. --- Hela \VUgras{mjölkkammaren} \textsuperscript{(59)} hålls mycket ren av
\VUgras{gårdsdrängen} \textsuperscript{(65)} som just har diskat
\VUgras{mjölkkrukorna} \textsuperscript{(62)} i den stora hink som han bär i
handen.

%%K t07-tit-9 sub
\VUsaut{3.0mm}
\VUpk{10.2pt}{40}{Fjäderfäna.}
\VUsaut{3.0mm}

%%K t07-c2-14-1 p
14. --- Med sina tjocka träskor går han litet tungt, och så skrämmer han bort
\VUgras{kalkonen} \textsuperscript{(68)}, \VUgras{ankorna}
\textsuperscript{(66)}, \VUgras{ankbondarna} \textsuperscript{(67)} och
\VUgras{gässen} \textsuperscript{(69)} som spärrar upp \VUgras{sin näbb}
\textsuperscript{(70)} och skriker oroligt.

%%K t07-c2-14-2 p
\VUgras{Tuppen} \textsuperscript{(71)}, mer stridslysten, reser sin röda
\VUgras{kam} \textsuperscript{(72)}, skakar fjädrarna i sin \VUgras{stjärt}
\textsuperscript{(73)} och låter höra ett ljudligt « kuckeliku ».

%%K t07-c2-15-1 p
15. --- En \VUgras{modershöna} \textsuperscript{(74)} letar på
\VUgras{gödselhögen} \textsuperscript{(81)} med sina \VUgras{små kycklingar}
\textsuperscript{(75)}. \VUgras{Grisen} \textsuperscript{(78)} flyr till sin mor, den enorma
\VUgras{suggan} \textsuperscript{(77)} som vi ser bredvid svinstian.

%%K t07-c2-16-1 p
16. --- I sina burar äter \VUgras{kaninerna} \textsuperscript{(79)} girigt det späda
\VUgras{gräset} \textsuperscript{(80)} som bondens dotter bär till dem. Lilla Johanna tycker
mycket om sina kaniner och, varje dag, kommer hon och besöker sina små vänner sedan hon har
hämtat de färska \VUgras{äggen} \textsuperscript{(76)} i hönshuset.
\VUsaut{6.0mm}
\VUfilet{22mm}
